% ===== 原文 pp.77--81 =====
% [原文 p.77]
注意，对 1-群上链 $\lambda$ 作替换
\begin{equation}
\beta\to\beta+\delta\lambda\,,
\tag{4.79}
\end{equation}
并不改变闭时空流形上的 $Z[A_1]$。于是我们得到如下陈述。

\begin{statementbox}{陈述 4.2：离散 $G^{(0)}$ 0-形式对称性的 2d SPT 相}
$d=2$ 时，离散 0-形式对称群 $G^{(0)}$ 的 SPT 相 $[\beta]$ 由系数取在 $\U(1)$ 中的第二群上同调的元素刻画
\begin{equation}
[\beta]\in H^2\left(G^{(0)},U(1)\right).
\tag{4.80}
\end{equation}
\end{statementbox}

% [Tutor Note: p.77 式 (4.80)——用 $H^2(G^{(0)},\U(1))$ 刻画 2d 玻色 SPT 相源于 Dijkgraaf--Witten：$\beta$ 是 $G^{(0)}$ 的 $\C^\times$ 值 2-上链，$\delta\beta=1$ 即上闭链条件（见式 (4.75)）。核对锚点：取 $G^{(0)}=\Z_2$，则 $H^2(\Z_2,\U(1))=\Z_2$，恰好对应 2d 的平凡相与一个非平凡相；若误写成 $H^1$，则 $H^1(\Z_2,\U(1))=\Z_2$ 分类的是 $G^{(0)}$ 的一维表示（荷），不是 SPT 相，故 2d SPT 的维数必须是 2。英文原词：group cohomology with coefficients in $\U(1)$（系数在 $\U(1)$ 中的群上同调）。]

我们可以容易地把 (4.76) 推广，得到

\begin{statementbox}{陈述 4.3：由 $G^{(0)}$ 保护的群上同调型 SPT 相}
离散 $G^{(0)}$ 0-形式对称性的 $d$ 维 SPT 相的一类由元素
\begin{equation}
[\alpha_d]\in H^d\left(G^{(0)},U(1)\right).
\tag{4.81}
\end{equation}
给出。该 SPT 相的配分函数可表示为
\begin{equation}
Z[A_1]=\int A_1^*\alpha_d\,.
\tag{4.82}
\end{equation}
\end{statementbox}

\subsubsection{反常理论与反常多项式}

为讨论反常理论，我们考虑 $(d+1)$ 维 SPT 相。

\para{开流形上的 SPT 相。}注意 SPT 相没有 't Hooft 反常，即对每个闭的 $(d+1)$ 维时空流形 $M_{d+1}$ 都有
\begin{equation}
Z[B_{p_i+1}+\delta\Lambda_{p_i}]=Z[B_{p_i+1}]\,,
\tag{4.83}
\end{equation}
然而，若我们考虑具有 $d$ 维边界的开时空 $M_{d+1}$ 上的 SPT 相
\begin{equation}
M_d=\partial M_{d+1}\,,
\tag{4.84}
\end{equation}
则 (4.83) 不再成立，而是有
\begin{equation}
\frac{Z[B_{p_i+1}+\delta\Lambda_{p_i}]}{Z[B_{p_i+1}]}=Z_\partial[B_{p_i+1},\Lambda_{p_i}]\,,
\tag{4.85}
\end{equation}

% [原文 p.78]
\begin{figure}[htbp]
\centering
\includegraphics[width=0.92\textwidth]{fig17.png}
\figtitle{17}{该图展示开 $(d+1)$ 维流形 $\mathcal{M}_{d+1}$ 上的 $(d+1)$ 维 SPT 相 $\mathcal{I}_{d+1}$，其边界为 $\mathcal{M}_d$；以及一个 $d$ 维理论 $\mathcal{T}$，其反常理论正是 $\mathcal{I}_{d+1}$。如右图所示把 $\mathcal{T}$ 沿 $\mathcal{M}_d$ 放置，可使复合系统无反常。}
\end{figure}

其中 $Z_\partial[B_{p_i+1},\Lambda_{p_i}]$ 完全局域在边界 $M_d$ 上，并且同时是背景场 $B_{p_i+1}$ 与沿 $M_d$ 的规范变换参数 $\Lambda_{p_i}$ 的函数。我们可以把 $Z_\partial$ 解释为在开流形上定义该 SPT 相时出现的 't Hooft 反常。

% [Tutor Note: p.78 式 (4.84)--(4.85)——$Z_\partial$ 只依赖 $B_{p_i+1}$ 与 $\Lambda_{p_i}$ 沿边界 $M_d$ 的取值，故它是 $(d+1)$ 维 SPT 相在带边界时的「反常流入」项。符号自洽性核对：定义 4.6 中给的是 $Z_\partial=\mathcal{I}_d^{-1}$（式 (4.88)），而下式 (4.86) 给的是 $\mathcal{I}_d=Z_\partial^{-1}$，二者互为逆，转录时不要漏掉负指数或把它改成 $+1$。]

\para{反常理论。}现在我们可以沿 $M_d$ 放置一个 $d$ 维理论 $\mathcal{T}$，它与该 SPT 相具有相同的对称性，但携带 't Hooft 反常
\begin{equation}
\mathcal{I}_d[B_{p_i+1},\Lambda_{p_i}]=Z_\partial^{-1}[B_{p_i+1},\Lambda_{p_i}]\,.
\tag{4.86}
\end{equation}
由 $(d+1)$ 维 SPT 与 $d$ 维理论 $\mathcal{T}$ 组成的混合 $(d,d+1)$ 维系统是无反常的。见图 17。

\begin{definitionbox}{定义 4.6：反常理论}
给定 $d$ 维理论的 't Hooft 反常
\begin{equation}
\mathcal{I}_d[B_{p_i+1},\Lambda_{p_i}]\,,
\tag{4.87}
\end{equation}
与之相联的反常理论 $Z[B_{p_i+1}]$ 是一个 $(d+1)$ 维 SPT 相，满足
\begin{equation}
Z_\partial[B_{p_i+1},\Lambda_{p_i}]=\mathcal{I}_d^{-1}[B_{p_i+1},\Lambda_{p_i}]\,.
\tag{4.88}
\end{equation}
\end{definitionbox}

在本文中，我们把反常理论记作 $\mathcal{I}_{d+1}$，其配分函数记作
\begin{equation}
\mathcal{I}_{d+1}[B_{p_i+1}]\,.
\tag{4.89}
\end{equation}

\para{反常多项式。}对连续对称性，按
\begin{equation}
\mathcal{I}_{d+1}[B_{p_i+1}]=\exp\left(2\pi\ii S^{\text{eff}}_{d+1}\right),
\tag{4.90}
\end{equation}
定义的反常理论 $\mathcal{I}_{d+1}$ 的有效作用量 $S^{\text{eff}}_{d+1}$ 不是闭的，而用它可定义一个 $(d+2)$ 维反常多项式
% [原文 p.79]
\begin{equation}
\mathcal{I}_{d+2}=\dd S^{\text{eff}}_{d+1}\,,
\tag{4.91}
\end{equation}
这与连续 0-形式对称性的研究中大家熟知的情形相同。

% [Tutor Note: p.79 式 (4.91)——下降（descent）程序的三层是：$\mathcal{I}_{d+1}=\exp(2\pi\ii S^{\text{eff}}_{d+1})$（式 (4.90)），$\mathcal{I}_{d+2}=\dd S^{\text{eff}}_{d+1}$（式 (4.91)），再往上 $\dd\mathcal{I}_{d+2}=0$（闭性）。自检锚点：对 Maxwell 理论 $S^{\text{eff}}_{d+1}=\int B_2^e\wedge\dd B^m_{d-2}$，故 $\mathcal{I}_{d+2}=\int \dd B_2^e\wedge \dd B^m_{d-2}$，正是式 (4.93)；若此处写成 $\mathcal{I}_{d+2}=\dd\mathcal{I}_{d+1}$ 就错了（维数差 1）。英文原词：descent procedure（下降程序）。]

't Hooft 反常 $\mathcal{I}_d$、反常理论 $\mathcal{I}_{d+1}$ 与反常多项式 $\mathcal{I}_{d+2}$ 通过下降程序相互联系。

\para{例子。}上文我们讨论了许多带 't Hooft 反常的例子。下面讨论相应的反常理论与反常多项式（如果存在的话）：

\begin{itemize}[leftmargin=1.4em,itemsep=1pt,topsep=1pt]
\item 对例 4.5 中讨论的 Maxwell 理论，反常理论为
\begin{equation}
\mathcal{I}_{d+1}=\exp\left(2\pi\ii\int B_2^e\wedge\dd B^m_{d-2}\right),
\tag{4.92}
\end{equation}
而反常多项式为
\begin{equation}
\mathcal{I}_{d+2}=\int \dd B_2^e\wedge \dd B^m_{d-2}\equiv\int F_3^e\wedge F^m_{d-1}\,,
\tag{4.93}
\end{equation}
其中 $F_3^e$ 与 $F^m_{d-1}$ 分别是电与磁背景场的场强。

\item 对例 4.6 中讨论的离散高阶形式规范理论，反常理论为
\begin{equation}
\mathcal{I}_{d+1}=\exp\left(-2\pi\ii s(p)\int B^e_{p+1}\cup_\eta B^m_{d-p}\right),\qquad s(p)=(-1)^p\,.
\tag{4.94}
\end{equation}

\item 对 (4.2.1) 中讨论的 2d 离散 0-形式对称群，反常理论为
\begin{equation}
\mathcal{I}_3=\int A_1^*\omega_3\,.
\tag{4.95}
\end{equation}
见式 (4.76) 之后的注记。
\end{itemize}

% [Tutor Note: p.79 式 (4.92)--(4.94)——注意三种书写方式的差别：连续对称性用微分形式与 $\wedge$（(4.92)、(4.93)），离散对称性用 cup 积 $\cup_\eta$（(4.94)）；$\cup_\eta$ 中的 $\eta$ 表示把 $\U(1)$ 值上链相乘时按 $\eta=\pm1$ 选取的配对（源于 cup 积的交换性符号）。次数核对：$B^e_2\wedge\dd B^m_{d-2}$ 的次数为 $2+(d-1)=d+1$，正是 $d+1$ 维反常理论的作用量；$\dd B_2^e\wedge\dd B^m_{d-2}$ 的次数为 $3+(d-1)=d+2$，正是 $d+2$ 维反常多项式。]

\subsubsection{非阿贝尔规范理论的混合 't Hooft 反常}

上一小节中，我们讨论了 $d$ 维非阿贝尔规范理论中电与磁高阶形式对称性的背景场。结果表明，与 Maxwell 情形一样，非阿贝尔规范理论在这两个对称性之间也存在混合 't Hooft 反常。

% [原文 p.80]
\para{三个障碍类。}回忆磁性对称性的背景场 $B^m_{d-2}$ 以作用量中的项 (4.27) 出现，其中注意 $w_2$ 是把 $\mathcal{G}$ 丛提升为 $G$ 丛的 $\mathcal{Z}$-值障碍类。另一方面，开启电对称性的背景场 $B_2^e$ 对应于考虑群 $\mathcal{G}/G^{(1)}$ 的丛，对这些丛而言，提升为 $\mathcal{G}$ 丛的 $G^{(1)}$-值障碍类就是背景场 $B_2^e$。

总体而言，我们可以直接考虑从
\begin{equation}
\frac{\mathcal{G}}{G^{(1)}}=\frac{G/\mathcal{Z}}{G^{(1)}}\equiv\frac{G}{\mathcal{X}}\,,\qquad \mathcal{X}\subseteq Z(G)\,,
\tag{4.96}
\end{equation}
丛到 $G$ 丛的提升，对后者我们将其 $\mathcal{X}$-值障碍类记为 $W_2$。如下文将要讨论的，结果表明：尽管 $W_2$ 与 $B_2^e$ 都是闭的，但既然 $G/\mathcal{X}$ 丛提升为 $G$ 丛、或提升为 $G/\mathcal{Z}$ 丛分别都有定义良好的障碍类，只要 $B_2^e$ 非平凡，场 $w_2$ 就不是闭的。这是因为 $w_2$ 是把 $G/\mathcal{Z}$ 丛提升为 $G$ 丛的障碍类，但当 $B_2^e$ 非零时，我们必须考虑那些一开始就不能视为 $G/\mathcal{Z}$ 丛的 $G/\mathcal{X}$ 丛。

% [Tutor Note: p.80 式 (4.96)——右端的等同用了 $\mathcal{G}=G/\mathcal{Z}$（即 $\mathcal{Z}$ 是把单连通 $G$ 商到 $\mathcal{G}$ 的中心子群）与式 (4.98) 的 $G^{(1)}=\mathcal{X}/\mathcal{Z}$：于是 $\mathcal{G}/G^{(1)}=(G/\mathcal{Z})/(\mathcal{X}/\mathcal{Z})=G/\mathcal{X}$。这也说明为什么要求 $\mathcal{X}\subseteq Z(G)$ 而不只是 $\mathcal{X}\subseteq\mathcal{G}$：只有 $\mathcal{X}$ 是 $G$ 的中心子群时，$\mathcal{X}/\mathcal{Z}$ 才自动正规，商 $\mathcal{G}/G^{(1)}$ 才是群而不是齐性空间。]

\para{联系各障碍类。}显然，$\mathcal{X}$ 应当是 $\mathcal{Z}$ 与 $G^{(1)}$ 的组合，而 $W_2$ 应当是 $w_2$ 与 $B_2^e$ 的组合。我们更详细地描述这些组合。首先，$\mathcal{Z}$ 自然是 $\mathcal{X}$ 的子群，从而给出一个单射同态
\begin{equation}
i:\ \mathcal{Z}\to\mathcal{X}\,.
\tag{4.97}
\end{equation}
把 $\mathcal{Z}$ 从 $\mathcal{X}$ 中扣除后剩下的部分是 $G^{(1)}$
\begin{equation}
G^{(1)}=\mathcal{X}/\mathcal{Z}\,,
\tag{4.98}
\end{equation}
从而给出一个满射同态
\begin{equation}
\pi:\ \mathcal{X}\to G^{(1)}\,.
\tag{4.99}
\end{equation}
把这两个同态结合起来，我们得到短正合列
\begin{equation}
0\to\mathcal{Z}\to\mathcal{X}\to G^{(1)}\to 0\,.
\tag{4.100}
\end{equation}
各障碍类之间的关系为
\begin{equation}
W_2=i(w_2)+\tilde{B}_2^e\,,
\tag{4.101}
\end{equation}
其中 $i(w_2)$ 是把同态 $i$ 作用到 $w_2$ 上得到的 $\mathcal{X}$-值 2-上链
\begin{equation}
i(w_2)(v_0,v_1,v_2)\equiv i\left(w_2(v_0,v_1,v_2)\in\mathcal{Z}\right)\in\mathcal{X}\,,
\tag{4.102}
\end{equation}
而 $\tilde{B}_2^e$ 是经由同态 $\pi$ 对 $B_2^e$ 所作的一个提升，即一个 $\mathcal{X}$-值 2-上链，满足
\begin{equation}
\pi(\tilde{B}_2^e)=B_2^e\,.
\tag{4.103}
\end{equation}

% [原文 p.81]
\para{障碍类 $w_2$ 的非闭性。}把微分 $\delta$ 作用到 (4.101) 两边，得到
\begin{equation}
0=i(\delta w_2)+\delta\tilde{B}_2^e\,,
\tag{4.104}
\end{equation}
因为 $\delta$ 与 $i$ 对易。现在，由于 $\delta$ 与 $\pi$ 对易，我们有
\begin{equation}
\pi(\delta\tilde{B}_2^e)=\delta B_2^e=0\,.
\tag{4.105}
\end{equation}
因此，我们可以通过把 $i$ 作用在一个 $\mathcal{Z}$-值 3-上链上来得到 $\delta\tilde{B}_2^e$，该 3-上链记作 $\mathrm{Bock}(B_2^e)$，称为 $B_2^e$ 关于短正合列 (4.100) 的 Bockstein
\begin{equation}
\delta\tilde{B}_2^e=i\left(\mathrm{Bock}(B_2^e)\right).
\tag{4.106}
\end{equation}
把它代回，我们得到关键关系
\begin{equation}
\delta w_2=-\mathrm{Bock}(B_2^e)\,,
\tag{4.107}
\end{equation}
它描述了 $w_2$ 的非闭性。

% [Tutor Note: p.81 式 (4.107)——负号不是约定选择，而是算出来的：由 (4.101) 且 $W_2$ 作为障碍类是闭的，有 $\delta W_2=i(\delta w_2)+\delta\tilde{B}_2^e=0$；代入 (4.106) 得 $i(\delta w_2)+i(\mathrm{Bock}(B_2^e))=i\big(\delta w_2+\mathrm{Bock}(B_2^e)\big)=0$，再用 $i$ 单射即得 $\delta w_2=-\mathrm{Bock}(B_2^e)$。核对锚点：取 $G^{(1)}=\Z_2,\ \mathcal{G}=\SO(3)$ 时 $\mathrm{Bock}$ 是 $\Z_2$ 的 Bockstein，$\delta w_2\neq0$ 与「$B_2^e$ 非零时 $w_2$ 不能整体定义」的几何图像一致。]

\begin{notebox}{上同调中的相伴长正合列}
给定 阿贝尔群的短正合列
\begin{equation}
0\to A\to B\to C\to 0\,,
\tag{4.108}
\end{equation}
我们得到上同调中的长正合列
\begin{equation}
\begin{gathered}
\cdots \longrightarrow H^p(M,A)\longrightarrow H^p(M,B)\longrightarrow H^p(M,C)\\
\xrightarrow{\ \mathrm{Bock}\ } H^{p+1}(M,A)\longrightarrow H^{p+1}(M,B)\longrightarrow H^{p+1}(M,C)\longrightarrow \cdots
\end{gathered}
\tag{4.109}
\end{equation}
其中 $M$ 是任意流形。每一行中的映射都直接来自短正合列作用在上链上的那些映射。两行之间的连接映射正是上文就 $p=2$ 讨论过的 Bockstein 映射。对一般的 $p$，Bockstein 映射是其直接推广，在上链上由
\begin{equation}
\delta\tilde{c}_p=i\left(\mathrm{Bock}(c_p)\right),
\tag{4.110}
\end{equation}
指定，其中 $c_p$ 是 $C$-值 $p$-上链。
\end{notebox}

% [Tutor Note: p.81 式 (4.109)——原文把长正合列画成两行（第一行到 $H^p(M,C)$、第二行从 $H^{p+1}(M,A)$ 起），并用一条弯曲箭头把 $H^p(M,C)$ 连到 $H^{p+1}(M,A)$、箭头上标 Bock；本译本改用两行居中排布，把该连接映射写为第二行行首的 $\xrightarrow{\mathrm{Bock}}$，数学内容不变。连接映射即 $\delta_*:H^p(M,C)\to H^{p+1}(M,A)$，它与 (4.106) 的关系是：对 $c_p$ 的类取 $\delta_*$ 后再用 $i$ 拉回，就等于上链层面的 $\delta\tilde c_p$。]

\para{计算 't Hooft 反常。}现在我们对磁性背景场 $B^m_{d-2}$ 作规范变换
\begin{equation}
B^m_{d-2}\to B^m_{d-2}+\delta\Lambda^m_{d-3}\,.
\tag{4.111}
\end{equation}
